\SourceLeaf{252}
$RR$ etc., etiam $Akk$ per $mm'$ divisibilem esse. Substitutis autem pro $a,2b,c$, $a',2b',c'$ valoribus suis $\frac{P}{n'}$ etc. sive $\frac{1}{n'}(pq'-qp')$ etc., transibunt in sex alias aequationes, in quibus ad dextram habebuntur producta ex quantitate $\frac{1}{nn'}(q'q''-qq''')$ in $PP,QQ,RR$ etc. Calculum facillimum lectoribus relinquimus. Hinc sequitur (quoniam omnes $PP,QQ$ etc. esse $=0$ nequeunt) $Ann'=q'q''-qq'''$.

Simili modo ex aequationibus fundamentalibus derivantur sex aliae aequationes, a praecedentibus in eo tantummodo discrepantes, quod pro $A$ ubique habetur $C$ et pro $q,q',q'',q'''$ resp. $p,p',p'',p'''$, quas ipsas brevitatis caussa non adscribimus. Hinc eodem modo sequitur, $Ckk$ per $mm'$ divisibilem esse atque $Cnn'=p'p''-pp'''$.

Denique ex eodem fonte petuntur sex aequationes hae:
\[
\begin{aligned}
BPP&=-aa'p'q'+ab'(pq'+qp')-ac'pq\\
BQQ&=-aa'p''q''+ba'(pq''+qp'')-ca'pq\\
BRR&=-aa'p'''q'''+(bb'+\Delta)(pq'''+qp''')-cc'pq\\
BSS&=-ac'p''q''+(bb'-\Delta)(p'q''+q'p'')-ca'p'q'\\
BTT&=-ac'p'''q'''+bc'(p'q'''+q'p''')-cc'p'q'\\
BUU&=-ca'p'''q'''+cb'(p''q'''+q''p''')-cc'p''q''
\end{aligned}
\]
unde perinde ut ante concluditur, $2Bkk$ divisibilem esse per $mm'$ atque $2Bnn'=pq'''+qp'''-p'q''-q'p''$.

Quoniam itaque $Akk,2Bkk,Ckk$ per $mm'$ sunt divisibiles, facile perspicietur, etiam $Mkk$ per $mm'$ divisibilem esse debere. Ex aequationibus fundamentalibus autem colligitur, $M$ metiri ipsos $aa',\allowbreak 2ab',\allowbreak ac',\allowbreak 2ba',\allowbreak 4bb',\allowbreak 2bc',\allowbreak ca',\allowbreak 2cb',\allowbreak cc'$, adeoque etiam ipsos $am',2bm',cm'$ (qui sunt divisores comm. max. trium primorum mediorum et ultimorum resp.); denique etiam ipsum $mm'$ qui est horum div. comm. max. Hinc patet, in eo casu ubi forma $F$ ex formis $f,f'$ composita est sive $k=1$, necessario esse $M=mm'$. Quae est \textsc{conclusio quinta}.

Si div. comm. max. numerorum $A,B,C$ est $\mathfrak{M}$, hic erit vel $=M$ (quando forma $F$ est proprie primitiva vel ex proprie primitiva derivata) vel $=\frac{1}{2}M$ (quando $F$ est forma improprie primitiva vel ex improprie prim. derivata); similiter designando divisores comm. max. numerorum $a,b,c;\ a',b',c'$ resp. per $\mathfrak{m},\mathfrak{m}'$ erit $\mathfrak{m}$ vel $=m$ vel $\frac{1}{2}m$, et $\mathfrak{m}'$ vel $=m'$ vel $=\frac{1}{2}m'$. Iam patet, $\mathfrak{m}\mathfrak{m}$ metiri ipsum $d'$, $\mathfrak{m}'\mathfrak{m}'$ ipsum $d'$, adeoque $\mathfrak{m}\mathfrak{m}\mathfrak{m}'\mathfrak{m}'$ ipsum $dd'$ sive $\Delta\Delta$, et $\mathfrak{m}\mathfrak{m}'$ ipsum $\Delta$.
\SourceLeaf{253}
Hinc ex sex ultimis aequationibus pro $BPP$ etc. sequitur, $\mathfrak{m}\mathfrak{m}'$ metiri ipsum $Bkk$, adeoque (quum etiam ipsos $Akk,Ckk$ metiatur) etiam ipsum $\mathfrak{M}kk$. Quoties igitur $F$ ex $f,f'$ composita est, metietur $\mathfrak{m}\mathfrak{m}'$ ipsum $\mathfrak{M}$. Quando itaque in hoc casu utraque $f,f'$ est proprie primitiva vel ex proprie primitiva derivata sive $\mathfrak{m}\mathfrak{m}'=mm'=M$, erit $\mathfrak{M}=M$, sive $F$ similis forma. Quando vero, in eadem suppositione, aut utraque $f,f'$ aut alterutra saltem est improprie primitiva vel ex improprie primitiva derivata, \textit{e. g.} forma $f$; ex aequationibus fundamentalibus sequitur, $aa',2ab',ac',ba',2bb',bc',ca',2cb',cc'$ per $\mathfrak{M}$ divisibiles esse adeoque etiam $am',bm',cm'$ et hinc quoque $\mathfrak{m}m'=\frac{1}{2}mm'=\frac{1}{2}M$; unde necessario in hoc casu erit $\mathfrak{M}=\frac{1}{2}M$, sive etiam forma $F$ vel impr. prim. vel ex impr. prim. derivata. Quae efficiunt \textsc{conclusionem sextam}.

Tandem observamus, si novem aequationes
\[
\begin{gathered}
an'=P,\quad 2bn'=R-S,\quad cn'=U\\
a'n=Q,\quad 2b'n=R+S,\quad c'n=T\\
\begin{gathered}Ann'=q'q''-qq''',\\2Bnn'=pq'''+qp'''-p'q''-q'p'',\\Cnn'=p'p''-pp'''\end{gathered}
\end{gathered}
\]
(quas, quoniam in sequentibus saepius ad ipsas revenire oportebit, per $\Omega$ designabimus) locum habere \textit{supponantur}, spectatis adeo ipsis $n,n'$ tamquam incognitis. quarum tamen neutra $=0$: per substitutionem facile confirmari, etiam aequationes fundamentales $1\ldots9$ necessario veras esse sive formam $(A,B,C)$ per substitutionem $p,p',p'',p''';\ q,q',q'',q'''$ in productum e formis $(a,b,c)(a',b',c')$ transire; praetereaque esse
\[
bb-ac=nn(BB-AC),\quad b'b'-a'c'=n'n'(BB-AC)
\]
Calculum quem hic apponere nimis prolixum foret lectorum industriae committimus.

\Needspace{8\baselineskip}
\Article{236}
\textsc{Problema.} \textit{Propositis duabus formis quarum determinantes aut aequales sunt aut saltem rationem quadratorum inter se habent: invenire formam ex illis compositam.}

\textit{Sol.} Sint formae componendae $(a,b,c)\ldots f$, $(a',b',c')\ldots f'$; harum determinantes $d,d'$; divisores communes maximi numerorum $a,2b,c;\ a',2b',c'$ resp. $m,m'$; divisor comm. maximus numerorum $dm'm',d'mm$ eodem signo ut $d,d'$
\SourceLeaf{254}
affectus $D$. Tunc $\frac{dm'm'}{D},\frac{d'mm}{D}$ erunt numeri positivi inter se primi ipsorumque productum, quadratum; quare ipsi erunt quadrata (art. 21). Hinc $\sqrt{\frac{d}{D}},\sqrt{\frac{d'}{D}}$ erunt quantitates rationales quas ponemus $=n,n'$, et quidem accipiemus pro $n$ valorem positivum vel negativum, prout forma $f$ in compositionem vel directe vel inverse ingredi debet, similiterque signum ipsius $n'$ ex ratione qua $f'$ in compositionem ingredi debet, determinabimus. Erunt itaque $mn',m'n$ numeri integri inter se primi; $n$ et $n'$ autem etiam fractiones esse possunt. His ita factis, observamus, $an',cn',a'n,c'n,bn'+b'n,bn'-b'n$ esse integros, quod de quatuor prioribus per se manifestum est (quum $an'=\frac{a}{m}mn'$ etc.); de duobus reliquis eodem modo probatur ut in art. praec. demonstratum fuit, $R$ et $S$ per $e$ divisibiles esse.

Iam accipiantur quatuor numeri integri $\mathfrak{Q},\mathfrak{Q}',\mathfrak{Q}'',\mathfrak{Q}'''$ ad libitum, ea sola conditione, ut quatuor quantitates in aequatione sequente (I) ad laevam positae non omnes simul $=0$ fiant, ponaturque
\[
\begin{aligned}
\mathfrak{Q}'an'+\mathfrak{Q}''a'n+\mathfrak{Q}'''(bn'+b'n)&=\mu q\qquad\ldotp\,\ldotp\,\ldotp\,\ldotp\,\ldotp\qquad\text{(I)}\\
-\mathfrak{Q}an'+\mathfrak{Q}'''c'n-\mathfrak{Q}''(bn'-b'n)&=\mu q'\\
\mathfrak{Q}'''cn'-\mathfrak{Q}a'n+\mathfrak{Q}'(bn'-b'n)&=\mu q''\\
-\mathfrak{Q}''cn'-\mathfrak{Q}'c'n-\mathfrak{Q}(bn'+b'n)&=\mu q'''
\end{aligned}
\]
ita ut $q,q',q'',q'''$ fiant integri divisorem communem non habentes, quod obtinetur accipiendo pro $\mu$ divisorem communem maximum quatuor numerorum, qui in his aequationibus sunt ad laevam. Tunc igitur per art. 40 inveniri poterunt quatuor numeri integri $\mathfrak{P},\mathfrak{P}',\mathfrak{P}'',\mathfrak{P}'''$ tales ut fiat
\[
\mathfrak{P}q+\mathfrak{P}'q'+\mathfrak{P}''q''+\mathfrak{P}'''q'''=1
\]
Quo facto determinentur numeri $p,p',p'',p'''$ per aequationes sequentes:
\[
\begin{aligned}
\mathfrak{P}'an'+\mathfrak{P}''a'n+\mathfrak{P}'''(bn'+b'n)&=p\qquad\ldotp\,\ldotp\,\ldotp\,\ldotp\,\ldotp\qquad\text{(II)}\\
-\mathfrak{P}an'+\mathfrak{P}'''c'n-\mathfrak{P}''(bn'-b'n)&=p'\\
\mathfrak{P}'''cn'-\mathfrak{P}a'n+\mathfrak{P}'(bn'-b'n)&=p''\\
-\mathfrak{P}''cn'-\mathfrak{P}'c'n-\mathfrak{P}(bn'+b'n)&=p'''
\end{aligned}
\]
Tandem ponatur
\[
\begin{gathered}q'q''-qq'''=Ann',\\pq'''+qp'''-p'q''-q'p''=2Bnn',\\p'p''-pp'''=Cnn'\end{gathered}
\]
\SourceLeaf{255}
Tunc $A,B,C$ erunt numeri integri formaque $(A,B,C)\ldots F$ ex formis $f,f'$ composita.

\textit{Dem.} I. Ex aequatt. I nullo negotio confirmantur sequentes quatuor aequationes:
\[
\begin{aligned}
0&=q'cn'-q''c'n-q'''(bn'-b'n)\qquad\ldotp\,\ldotp\,\ldotp\,\ldotp\,\ldotp\,\ldotp\qquad\text{(III)}\\
0&=qcn'+q'''a'n-q''(bn'+b'n)\\
0&=q'''an'+qc'n-q'(bn'+b'n)\\
0&=q''an'-q'a'n-q(bn'-b'n)
\end{aligned}
\]

II. Iam ponamus, numeros integros $\mathfrak{A},\mathfrak{B},\mathfrak{C},\mathfrak{A}',\mathfrak{B}',\mathfrak{C}',\mathfrak{N},\mathfrak{N}'$ ita determinatos esse ut fiat
\[
\begin{aligned}
\mathfrak{A}a+2\mathfrak{B}b+\mathfrak{C}c&=m\\
\mathfrak{A}'a'+2\mathfrak{B}'b'+\mathfrak{C}'c'&=m'\\
\mathfrak{N}m'n+\mathfrak{N}'mn'&=1
\end{aligned}
\]
Tunc erit
\[
\begin{aligned}\mathfrak{A}a\mathfrak{N}'n'&+2\mathfrak{B}b\mathfrak{N}'n'+\mathfrak{C}c\mathfrak{N}'n'\\&+\mathfrak{A}'a'\mathfrak{N}n+2\mathfrak{B}'b'\mathfrak{N}n+\mathfrak{C}'c'\mathfrak{N}n=1\end{aligned}
\]
Hinc atque ex aequatt. (III) facile confirmatur, si statuatur
\[
\begin{aligned}
-q'\mathfrak{A}\mathfrak{N}'-q''\mathfrak{A}'\mathfrak{N}-q'''(\mathfrak{B}\mathfrak{N}'+\mathfrak{B}'\mathfrak{N})&=\mathfrak{q}\\
q\mathfrak{A}\mathfrak{N}'-q'''\mathfrak{C}'\mathfrak{N}+q''(\mathfrak{B}\mathfrak{N}'-\mathfrak{B}'\mathfrak{N})&=\mathfrak{q}'\\
-q'''\mathfrak{C}\mathfrak{N}'+q\mathfrak{A}'\mathfrak{N}-q'(\mathfrak{B}\mathfrak{N}'-\mathfrak{B}'\mathfrak{N})&=\mathfrak{q}''\\
q''\mathfrak{C}\mathfrak{N}'+q'\mathfrak{C}'\mathfrak{N}+q(\mathfrak{B}\mathfrak{N}'+\mathfrak{B}'\mathfrak{N})&=\mathfrak{q}'''
\end{aligned}
\]
fore
\[
\begin{aligned}
\mathfrak{q}'an'+\mathfrak{q}''a'n+\mathfrak{q}'''(bn'+b'n)&=q\qquad\ldotp\,\ldotp\,\ldotp\,\ldotp\,\ldotp\,\ldotp\qquad\text{(IV)}\\
-\mathfrak{q}an'+\mathfrak{q}'''c'n-\mathfrak{q}''(bn'-b'n)&=q'\\
\mathfrak{q}'''cn'-\mathfrak{q}a'n+\mathfrak{q}'(bn'-b'n)&=q''\\
-\mathfrak{q}''cn'-\mathfrak{q}'c'n-\mathfrak{q}(bn'+b'n)&=q'''
\end{aligned}
\]

Quoties $\mu=1$, hae aequationes non sunt necessariae, sed ipsarum loco aequationes (I), quibus omnino analogae sunt, retineri possunt. Quodsi nunc ex aequatt. II, IV valores ipsorum $Ann',2Bnn',Cnn'$ (\textit{i. e.} numerorum $q'q''-qq'''$ etc.) evolvuntur, et quae mutuo se destruunt delentur: invenietur, singulorum
\SourceLeaf{256}
partes esse vel producta ex integris in $nn'$, vel ex integris in $dn'n'$ vel ex integris in $d'nn$, insuperque omnes partes constituentes ipsius $2Bnn'$ implicare factorem 2. Hinc concluditur (quoniam $dn'n'=d'nn$, et proin $\frac{dn'n'}{nn'}=\frac{d'nn}{nn'}=\sqrt{dd'}$ sunt integri), $A,B,C$ esse numeros integros. \textit{Q. E. P.}

III. Substituendo ex aequatt. (II) valores ipsorum $p,p',p'',p'''$, facile comprobatur adiumento aequatt. (III) et huius
\[
\mathfrak{P}q+\mathfrak{P}'q'+\mathfrak{P}''q''+\mathfrak{P}'''q'''=1
\]
esse
\[
\begin{aligned}
pq'-qp'&=an',\\ pq'''-qp'''-p'q''+q'p''&=2bn',\\p''q'''-q''p'''&=cn'\\
pq''-qp''&=a'n,\\ pq'''-qp'''+p'q''-q'p''&=2b'n,\\p'q'''-q'p'''&=c'n
\end{aligned}
\]
quae aequationes identicae sunt cum sex prioribus ($\Omega$) art. praec.; tres reliquae autem iam per hyp. locum habent. Quare (\textit{ibid. sub fin.}) forma $F$ transibit in $ff'$ per substitutionem $p,p',p'',p''';\ q,q',q''\,q'''$; ipsiusque determinans erit $=D$, sive aequalis divis. comm. max. numerorum $dm'm',d'mm$, quamobrem per concl. quartam art. praec. $F$ ex $f,f'$ composita erit. \textit{Q. E. S.} Denique facile perspicietur, $F$ ex $f,f'$ \textit{ita} compositam esse ut praescriptum sit, quum signa quantitatum $n,n'$ iam ab initio rite sint determinata.

\Needspace{8\baselineskip}
\Article{237}
\textsc{Theorema.} \textit{Si forma $F$ in productum e duabus formis $f,f'$ est transformabilis, atque forma $f'$ formam $f''$ implicat: $F$ etiam in productum e formis $f,f''$ transformabilis erit.}

\textit{Dem.} Retineantur pro formis $F,f,f'$ omnia signa art. 235; forma $f''$ sit $=(a'',b'',c'')$, transeatque $f'$ in $f''$ per substitutionem $\alpha.\ \beta,\gamma,\delta$. Tunc nullo negotio perspicietur, $F$ transire in $ff''$ per substitutionem
\[
\begin{array}{llll}
\alpha p+\gamma p',&\beta p+\delta p',&\alpha p''+\gamma p''',&\beta p''+\delta p'''\\
\alpha q+\gamma q',&\beta q+\delta q',&\alpha q''+\gamma q''',&\beta q''+\delta q'''
\end{array}
\qquad\text{\textit{Q. E. D.}}
\]

Positis brevitatis caussa coefficientibus
\[
\begin{gathered}\alpha p+\gamma p',\quad\beta p+\delta p'\ \text{etc.}\\=\mathfrak{P},\mathfrak{P}',\mathfrak{P}'',\mathfrak{P}''';\quad\mathfrak{Q},\mathfrak{Q}',\mathfrak{Q}'',\mathfrak{Q}'''\end{gathered}
\]
\SourceLeaf{257}
numeroque $\alpha\delta-\beta\gamma=e$: ex aequatt. $\Omega$ art. 235 facile confirmatur, esse
\[
\begin{aligned}
\mathfrak{P}\mathfrak{Q}'-\mathfrak{Q}\mathfrak{P}'&=an'e\\
\mathfrak{P}\mathfrak{Q}'''-\mathfrak{Q}\mathfrak{P}'''-\mathfrak{P}'\mathfrak{Q}''+\mathfrak{Q}'\mathfrak{P}''&=2bn'e\\
\mathfrak{P}''\mathfrak{Q}'''-\mathfrak{Q}''\mathfrak{P}'''&=cn'e\\
\mathfrak{P}\mathfrak{Q}''-\mathfrak{Q}\mathfrak{P}''&=\alpha\alpha a'n+2\alpha\gamma b'n+\gamma\gamma c'n=a''n\\
\mathfrak{P}\mathfrak{Q}'''-\mathfrak{Q}\mathfrak{P}'''+\mathfrak{P}'\mathfrak{Q}''-\mathfrak{Q}'\mathfrak{P}''&=2b''n\\
\mathfrak{P}'\mathfrak{Q}'''-\mathfrak{Q}'\mathfrak{P}'''&=c''n\\
\mathfrak{Q}'\mathfrak{Q}''-\mathfrak{Q}\mathfrak{Q}'''&=Ann'e\\
\mathfrak{P}\mathfrak{Q}'''+\mathfrak{Q}\mathfrak{P}'''-\mathfrak{P}'\mathfrak{Q}''-\mathfrak{Q}'\mathfrak{P}''&=2Bnn'e\\
\mathfrak{P}'\mathfrak{P}''-\mathfrak{P}\mathfrak{P}'''&=Cnn'e
\end{aligned}
\]
Iam designato determinante formae $f''$ per $d''$, erit $e$ radix quadrata ex $\frac{d''}{d'}$, et quidem positiva vel negativa, prout forma $f'$ formam $f''$ vel proprie vel improprie implicat. Quare $n'e$ erit radix quadrata ex $\frac{d''}{D}$; unde patet, novem aequationes praecedentes aequationibus $\Omega$ art. 235 prorsus analogas esse, formamque $f$ in transformatione formae $F$ in $ff''$ eodem modo accipi, ut in transformatione formae $F$ in $ff'$; formam $f''$ vero in illa vel eodem modo ut $f'$ in hac, vel opposito, prout $f'$ ipsam $f''$ proprie implicet vel improprie.

\Needspace{8\baselineskip}
\Article{238}
\textsc{Theorema.} \textit{Si forma $F$ sub forma $F'$ est contenta atque in productum e formis $f,f'$ transformabilis: etiam forma $F'$ in idem productum transformabilis erit.}

\textit{Dem.} Retentis pro formis $F,f,f'$ iisdem signis ut supra et supponendo formam $F'$ transire in $F$ per substitutionem $\alpha,\beta,\gamma,\delta$, facile perspicietur, $F'$ per substitutionem
\[
\begin{array}{llll}
\alpha p+\beta q,&\alpha p'+\beta q',&\alpha p''+\beta q'',&\alpha p'''+\beta q'''\\
\gamma p+\delta q,&\gamma p'+\delta q',&\gamma p''+\delta q'',&\gamma p'''+\delta q'''
\end{array}
\]
idem fieri quod $F$ per substitutionem $p,p',p'',p''';\ q,q',q'',q'''$, adeoque $F'$ per substitutionem illam transire in $ff'$. \textit{Q. E. D.}

Praeterea per similem calculum ut in art. praec. facile confirmatur, $F'$ eodem modo in $ff'$ transformabilem fore ut $F$, quando $F'$ ipsam $F$ proprie implicet; quando vero $F$ improprie sub $F'$ contenta sit, transformationes formae $F$ in $ff'$ et formae $F'$ in $ff'$ oppositas fore respectu utriusque formae $f,f'$, scilicet
\SourceLeaf{258}
quae ex his formis in alteram transformationem directe ingrediatur, in altera accipi inverse.

Ex combinatione theorematis praesentis cum theor. art. praec. obtinemus sequens generalius: \textit{Si forma $F$ in productum $ff'$ est transformabilis, atque formae $f,f'$ resp. implicant formas $g,g'$, forma $F$ vero sub forma $G$ contenta est: $G$ in productum $gg'$ transformabilis erit.} Nam per theor. art. praes. $G$ transformabilis erit in $ff'$, hinc per theor. art. praec. in $fg'$ et per idem theor. etiam in $gg'$. Porro patet, si omnes tres formae $f,f',G$ formas $g,g',F$ proprie implicent, $G$ eodem modo in $gg'$ transformabilem fore respectu formarum $g,g'$, ut $F$ in $ff'$ respectu formarum $f,f'$; idem evenire, si illae tres implicationes omnes sint impropriae; denique aeque facile determinari poterit, \textit{quomodo} $G$ in $gg'$ transformabilis sit, si ex illis implicationibus aliqua duabus reliquis sit dissimilis.

Si formae $F,f,f'$ formis $G,g,g'$ resp. sunt aequivalentes, hae eosdem determinantes habebunt ut illae, et quod pro formis $f,f'$ sunt numeri $m,m'$, idem erunt pro formis $g,g'$ (art. 161). Hinc nullo negotio per conclus. quartam art. 235 deducitur, in hocce casu $G$ ex $g,g'$ \textit{compositam} fore, si $F$ ex $f,f'$ composita sit, et quidem formam $g$ in compositionem illam eodem modo ingredi, ut $f$ in hac, quando $F$ ipsi $G$ eodem modo aequivaleat, ut $f$ ipsi $g$, et contra; similiterque $g'$ in compositione priori vel eodem modo vel opposito accipiendam ut $f'$ in posteriori, prout aequivalentia formarum $f',g'$ aequivalentiae formarum $F,G$ similis sit vel dissimilis.

\Needspace{8\baselineskip}
\Article{239}
\textsc{Theorema.} \textit{Si forma $F$ ex formis $f,f'$ composita est: quaevis alia forma in productum $ff'$ eodem modo transformabilis ut $F$, ipsam $F$ proprie implicabit.}

\textit{Dem.} Retentis pro $F,f,f'$ omnibus signis art. 235, aequationes $\Omega$ etiam hic locum habebunt. Ponamus formam $F'=(A',B',C')$, cuius determinans $=D'$, transire in productum $ff'$ per substitutionem $\mathfrak{p},\mathfrak{p}',\mathfrak{p}'',\mathfrak{p}''';\ \mathfrak{q},\mathfrak{q}',\mathfrak{q}'',\mathfrak{q}'''$ designemusque numeros
\[
\begin{gathered}\mathfrak{p}\mathfrak{q}'-\mathfrak{q}\mathfrak{p}',\quad
\mathfrak{p}\mathfrak{q}''-\mathfrak{q}\mathfrak{p}'',\quad
\mathfrak{p}\mathfrak{q}'''-\mathfrak{q}\mathfrak{p}'''.\\
\mathfrak{p}'\mathfrak{q}''-\mathfrak{q}'\mathfrak{p}'',\quad
\mathfrak{p}'\mathfrak{q}'''-\mathfrak{q}'\mathfrak{p}''',\quad
\mathfrak{p}''\mathfrak{q}'''-\mathfrak{q}''\mathfrak{p}'''\end{gathered}
\]
resp. per
\[
P',Q',R',S',T',U'
\]
\SourceLeaf{259}
Tunc habebuntur novem aequationes ipsis $\Omega$ omnino similes puta
\[
\begin{gathered}
P'=a\mathfrak{n}',\quad R'-S'=2b\mathfrak{n}',\quad U'=c\mathfrak{n}'\\
Q'=a'\mathfrak{n},\quad R'+S'=2b'\mathfrak{n},\quad T'=c'\mathfrak{n}\\
\mathfrak{q}'\mathfrak{q}''-\mathfrak{q}\mathfrak{q}'''=A'\mathfrak{n}\mathfrak{n}',\\
\mathfrak{p}\mathfrak{q}'''+\mathfrak{q}\mathfrak{p}'''-\mathfrak{p}'\mathfrak{q}''-\mathfrak{q}'\mathfrak{p}''=2B'\mathfrak{n}\mathfrak{n}',\\
\mathfrak{p}'\mathfrak{p}''-\mathfrak{p}\mathfrak{p}'''=C'\mathfrak{n}\mathfrak{n}'
\end{gathered}
\]
quas per $\Omega'$ designabimus. Quantitates $\mathfrak{n},\mathfrak{n}'$ hic erunt radices quadratae ex $\frac{d}{D'},\frac{d'}{D'}$ et quidem iisdem signis resp. affectae ut $n,n'$; si igitur radix quadrata ex $\frac{D}{D'}$ positive accepta (quae erit numerus integer) statuitur $=k$, erit $\mathfrak{n}=kn$, $\mathfrak{n}'=kn'$. Hinc et ex aequatt. senis prioribus in $\Omega$ et $\Omega'$ manifestum est, fore
\[
\begin{gathered}
P'=kP,\quad Q'=kQ,\quad R'=kR\\
S'=kS,\quad T'=kT,\quad U'=kU
\end{gathered}
\]
Quare per lemma art. 234 determinari poterunt quatuor numeri integri $\alpha,\beta,\gamma,\delta$ tales ut fiat
\[
\begin{aligned}
\alpha p+\beta q&=\mathfrak{p}.&\gamma p+\delta q&=\mathfrak{q}\\
\alpha p'+\beta q'&=\mathfrak{p}',&\gamma p'+\delta q'&=\mathfrak{q}'\ \text{etc.}
\end{aligned}
\]
atque
\[
\alpha\delta-\beta\gamma=k
\]
Substitutis his valoribus ipsorum $\mathfrak{p},\mathfrak{q},\mathfrak{p}',\mathfrak{q}'$ etc. in aequatt. tribus ultimis $\Omega'$, facile confirmatur adiumento aequationum $\mathfrak{n}=kn,\ \mathfrak{n}'=kn'$ triumque ultimarum $\Omega$, fore
\[
\begin{aligned}
A'\alpha\alpha+2B'\alpha\gamma+C'\gamma\gamma&=A\\
A'\alpha\beta+B'(\alpha\delta+\beta\gamma)+C'\gamma\delta&=B\\
A'\beta\beta+2B'\beta\delta+C'\delta\delta&=C
\end{aligned}
\]
quapropter forma $F'$ per substitutionem $\alpha,\beta,\gamma,\delta$ (quae propria erit, quoniam $\alpha\delta-\beta\gamma=k$ est positivus) transibit in $F$, \textit{i. e.} formam $F$ proprie implicabit. \textit{Q. E. D.}

Si itaque $F'$ e formis $f,f'$ etiam composita est (eodem modo ut $F$ ex iisdem), formae $F,F'$ eundem determinantem habebunt, eruntque adeo proprie aequivalentes. Generalius, si forma $G$ e formis $g,g'$ eodem modo composita est ut
\SourceLeaf{260}
$F$ ex $f,f'$ resp., formaeque $g,g'$ ipsis $f,f'$ proprie aequivalent: formae $F,G$ proprie aequivalebunt.

Quum is casus ubi ambae formae componendae compositionem directe ingrediuntur, simplicissimus sit, ad ipsumque reliqui facile reducantur, illum solum in sequentibus contemplabimur, ita ut si forma aliqua simpliciter dicatur e duabus aliis composita, semper subintelligere oporteat, ex utraque illam proprie esse compositam\HistoricalNote{*)}{Similiter ut in compositione rationum (quae cum compositione formarum magnam analogiam habet) subintelligi solet, rationes componendas directe accipiendas esse nisi ubi contrarium monetur.}. Eadem restrictio valebit, quoties forma in productum e duabus aliis transformabilis dicetur.

\Needspace{8\baselineskip}
\Article{240}
\textsc{Theorema.} \textit{Si e formis $f,f'$ composita est forma $F$; ex $F$ et $f''$ forma $\mathfrak{F}$; ex $f,f''$ forma $F'$; ex $F'$ et $f'$ forma $\mathfrak{F}'$: formae $\mathfrak{F},\mathfrak{F}'$ proprie aequivalentes erunt.}

\textit{Dem.} I. Sit
\[
\begin{aligned}
f&=axx+2bxy+cyy\\
f'&=a'x'x'+2b'x'y'+c'y'y'\\
f''&=a''x''x''+2b''x''y''+c''y''y''\\
F&=AXX+2BXY+CYY\\
F'&=A'X'X'+2B'X'Y'+C'Y'Y'\\
\mathfrak{F}&=\mathfrak{A}\mathfrak{X}\mathfrak{X}+2\mathfrak{B}\mathfrak{X}\mathfrak{Y}+\mathfrak{C}\mathfrak{Y}\mathfrak{Y}\\
\mathfrak{F}'&=\mathfrak{A}'\mathfrak{X}'\mathfrak{X}'+2\mathfrak{B}'\mathfrak{X}'\mathfrak{Y}'+\mathfrak{C}'\mathfrak{Y}'\mathfrak{Y}'
\end{aligned}
\]
determinantes harum septem formarum resp. $d,d',d'',D,D',\mathfrak{D},\mathfrak{D}'$, qui omnes eadem signa et rationem quadratorum inter se habebunt. Porro sit $m$ divisor communis maximus numerorum $a,2b,c$, similemque significationem habeant $m',m'',M$ relative ad formas $f',f'';F$. Tum ex concl. 4 art. 235, $D$ erit div. comm. max. numerorum $dm'm',d'mm$ adeoque $Dm''m''$ div. comm. max. numerorum $dm'm'm''m'',d'mmm''m''$; $M=mm'$; $\mathfrak{D}$ div. comm. max. num. $Dm''m''$, $d''MM$, sive numerorum $Dm''m'',d''mmm'm'$. Hinc concluditur, $\mathfrak{D}$ esse div. comm. max. trium numerorum $dm'm'm''m'',d'mmm''m'',d''mmm'm'$; ex simili autem ratione $\mathfrak{D}'$ eorundem trium numerorum divisor communis maximus erit; quare quum $\mathfrak{D},\mathfrak{D}'$ eadem signa habeant, erit $\mathfrak{D}=\mathfrak{D}'$, sive formae $\mathfrak{F},\mathfrak{F}'$ eundem determinantem habebunt.

\SourceLeaf{261}
II. Iam transeat $F$ in $ff'$ per substitutionem
\[
\begin{aligned}
X&=pxx'+p'xy'+p''yx'+p'''yy'\\
Y&=qxx'+q'xy'+q''yx'+q'''yy'
\end{aligned}
\]
atque $\mathfrak{F}$ in $Ff''$ per substitutionem
\[
\begin{aligned}
\mathfrak{X}&=\mathfrak{p}Xx''+\mathfrak{p}'Xy''+\mathfrak{p}''Yx''+\mathfrak{p}'''Yy''\\
\mathfrak{Y}&=\mathfrak{q}Xx''+\mathfrak{q}'Xy''+\mathfrak{q}''Yx''+\mathfrak{q}'''Yy''
\end{aligned}
\]
designenturque radices quadratae positivae ex $\frac{d}{D},\frac{d'}{D},\frac{D}{\mathfrak{D}}.\ \frac{d''}{\mathfrak{D}}$ per $n,n',\mathfrak{N},\mathfrak{n}''$. Tunc per art. 235 habebuntur decem et octo aequationes, quarum semissis altera ad transformationem formae $F$ in $ff'$ pertinebit, altera ad transformationem formae $\mathfrak{F}$ in $Ff''$. Prima erit $pq'-qp'=an'$, ad cuius instar facile formari poterunt reliquae brevitatis gratia hic omittendae. Ceterum quantitates $n,n',\mathfrak{N},\mathfrak{n}''$ rationales quidem erunt, sed non necessario numeri integri.

III. Si valores ipsorum $X,Y$ in valoribus ipsorum $\mathfrak{X},\mathfrak{Y}$ substituuntur. prodit substitutio talis:
\[
\begin{aligned}
\mathfrak{X}={}&(1)xx'x''+(2)xx'y''+(3)xy'x''+(4)xy'y''\\
&+(5)yx'x''+(6)yx'y''+(7)yy'x''+(8)yy'y''\\
\mathfrak{Y}={}&(9)xx'x''+(10)xx'y''+(11)xy'x''+(12)xy'y''\\
&+(13)yx'x''+(14)yx'y''+(15)yy'x''+(16)yy'y''
\end{aligned}
\]
per quam manifesto $\mathfrak{F}$ transibit in productum $ff'f''$. Coefficiens $(1)$ erit $=p\mathfrak{p}+q\mathfrak{p}''$; valores quindecim reliquorum non apponimus, quippe quos quisque nullo negotio evolvet. Designemus numerum $(1)(10)-(2)(9)$ per $(1,2)$, numerum $(1)(11)-(3)(9)$ per $(1,3)$, et generaliter $(g)(8+h)-(h)(8+g)$ per $(g,h)$, supponendo $g,h$ esse integros inaequales inter 1 et 16, quorum maior $h$\HistoricalNote{*)}{Horum signorum significatio praesens non est confundenda cum ea, in qua in art. 234 accepta erant; nam numeri per haec signa \textit{hic} expressi apprime respondent iis, qui in art. 234 per numeros similibus signis \textit{illic} denotatos multiplicabantur.}; hoc modo omnino viginti et octo signa habebuntur. Iam denotatis radicibus quadratis positivis ex $\frac{d}{\mathfrak{D}},\frac{d'}{\mathfrak{D}}$ per $\mathfrak{n},\mathfrak{n}'$. (quae erunt $=n\mathfrak{N},n'\mathfrak{N}$). eruentur sequentes 28 aequationes:
